\notechapter{N-20260302}{Pure Mathematics as a Language for Machine Learning and Data Analysis}

An AI-assisted attempt to map pure-mathematical languages onto machine learning kept returning to the same phenomenon: coordinates are useful until the data begin to carry geometry, symmetry, topology, or distributions of their own.  Backpropagation is the first place where this structural language becomes completely concrete.

\section{Backpropagation as cotangent pullback}

The previous backpropagation calculation can be read in a categorical way without making it mystical.  A forward neural network is a composite of maps between parameter spaces, activation spaces, and finally the one-dimensional loss space.  If
\[
  L:Y\to \mathbb R
\]
is the loss, then its differential at the output is a cotangent vector
\[
  dL_y\in T_y^*Y.
\]
Backpropagation pulls this covector backward along the forward computational graph.  In this sense the derivative operation used in training is contravariant:
\[
  f:X\to Y
  \quad\Longrightarrow\quad
  f^*:T^*Y\to T^*X.
\]
Forward maps push tangent perturbations forward; reverse-mode automatic differentiation pulls covectors backward.

For a concrete PyTorch-shaped layer, take
\[
  X\in\mathbb R^{b\times d},\qquad
  W\in\mathbb R^{n\times d},\qquad
  Y=XW^T\in\mathbb R^{b\times n}.
\]
In indices,
\[
  Y_{\alpha i}=\sum_{k=1}^d X_{\alpha k}W_{ik}.
\]
The full Jacobian with respect to the weight matrix has entries
\[
  \frac{\partial Y_{\alpha i}}{\partial W_{jk}}
  =\delta_{ij}X_{\alpha k},
\]
so its tensor shape is
\[
  (b,n,n,d).
\]
This is the honest tensor object.  Neural-network code does not materialize it because the upstream gradient
\[
  G=\frac{\partial L}{\partial Y}\in\mathbb R^{b\times n}
\]
immediately contracts it:
\[
\begin{aligned}
  \frac{\partial L}{\partial W_{jk}}
  &=
  \sum_{\alpha=1}^b\sum_{i=1}^n
  \frac{\partial L}{\partial Y_{\alpha i}}
  \frac{\partial Y_{\alpha i}}{\partial W_{jk}}\\
  &=
  \sum_{\alpha=1}^b G_{\alpha j}X_{\alpha k}.
\end{aligned}
\]
Therefore
\[
  \boxed{\displaystyle \frac{\partial L}{\partial W}=G^TX.}
\]
If one stores weights in the column convention \(W_{\mathrm{col}}\in\mathbb R^{d\times n}\) and writes \(Y=XW_{\mathrm{col}}\), the same contraction is
\[
  \frac{\partial L}{\partial W_{\mathrm{col}}}=X^TG.
\]
The categorical phrase ``pull back the covector'' is exactly this matrix multiplication: the high-order Jacobian collapses because it is being paired with a covector from the scalar loss.

\section{Geometric data analysis}

The same question reappears in data analysis: which mathematical language becomes necessary when coordinates alone no longer describe the structure?

A point cloud is often treated as a finite sample from a hidden space.  Topological data analysis makes this idea precise by replacing the cloud with a filtered simplicial complex.  For a scale parameter $\varepsilon$, one forms for instance the Vietoris--Rips complex
\[
  VR_\varepsilon(X)=\{\sigma\subseteq X:\ d(x,y)\le \varepsilon\text{ for all }x,y\in \sigma\}.
\]
As $\varepsilon$ grows, connected components merge, loops appear and disappear, and higher-dimensional holes may be born and die.  Persistent homology records this evolution.  The chain complex viewpoint is governed by the boundary identity
\[
  \partial_{k-1}\circ \partial_k=0,
\]
and homology is
\[
  H_k=\ker \partial_k/\operatorname{im}\partial_{k+1}.
\]
This quotient is the mathematical form of the phrase ``holes that are cycles but not boundaries.''

\section{Manifold hypothesis and embeddings}

The manifold hypothesis says that high-dimensional data often concentrate near a lower-dimensional manifold $M\subseteq \mathbb R^N$.  If $M$ is a smooth $d$-manifold, Whitney's embedding theorem says that $M$ can be embedded in $\mathbb R^{2d}$, and generically in a slightly larger ambient space.  In machine learning this gives a conceptual reason why dimensionality reduction is not only compression: it tries to recover the intrinsic coordinates of $M$.

The practical warning is that an embedding preserves the manifold as a smooth object, but not necessarily distances, probability densities, or task-relevant semantics.  PCA preserves a linear subspace model; manifold learning tries to preserve local neighborhoods; representation learning tries to learn coordinates useful for a loss function.

\section{Optimal transport and stochastic dynamics}

For probability measures $\mu,\nu$ on a metric space, the $p$-Wasserstein distance is
\[
  W_p(\mu,\nu)=\left(\inf_{\pi\in \Pi(\mu,\nu)}\int d(x,y)^p\,d\pi(x,y)\right)^{1/p}.
\]
Here $\Pi(\mu,\nu)$ is the set of couplings with marginals $\mu$ and $\nu$.  The dual viewpoint is equally important.  For $p=1$,
\[
  W_1(\mu,\nu)=\sup_{\|f\|_{\mathrm{Lip}}\le1}\left(\int f\,d\mu-\int f\,d\nu\right).
\]

For \(p=2\) there is a useful one-dimensional closed form.  Let
\[
  \mu=\mathcal N(a,\sigma_1^2),
  \qquad
  \nu=\mathcal N(b,\sigma_2^2),
  \qquad \sigma_1,\sigma_2>0.
\]
In one dimension the optimal map between these two Gaussians is affine and monotone, so write
\[
  T(x)=\alpha x+\beta.
\]
The condition \(T_\#\mu=\nu\) forces
\[
  \alpha=\frac{\sigma_2}{\sigma_1},
  \qquad
  \beta=b-\alpha a.
\]
If \(X=a+\sigma_1 Z\) with \(Z\sim\mathcal N(0,1)\), then
\[
  T(X)=b+\sigma_2 Z.
\]
Thus
\[
\begin{aligned}
  W_2^2(\mu,\nu)
  &=\mathbb E\bigl[(X-T(X))^2\bigr]\\
  &=\mathbb E\bigl[((a-b)+(\sigma_1-\sigma_2)Z)^2\bigr]\\
  &=(a-b)^2+(\sigma_1-\sigma_2)^2.
\end{aligned}
\]
This formula is small, but it explains a lot of geometric language around generative models: matching Gaussians means matching both their centers and their spreads.

Diffusion models can be read through stochastic differential equations
\[
  dX_t=b(t,X_t)\,dt+\sigma(t)\,dW_t.
\]
The density $p_t$ evolves according to a Fokker--Planck equation, and the reverse-time dynamics involve the score $\nabla_x\log p_t(x)$.  Thus score-based generative modeling is an attempt to learn the vector field that reverses a noisy transport of probability mass.

For DDPMs the training objective is usually written as a variational lower bound, or in its simplified form as a denoising mean-square loss.  It is not literally a Wasserstein objective.  Still, each reverse step compares Gaussian conditional distributions, and for Gaussians the geometry above says that a squared displacement of means and a mismatch of variances are exactly the ingredients of \(W_2^2\).  The ELBO uses KL divergences between Gaussian conditionals; when the variance is fixed or separately prescribed, minimizing that KL has the same practical center-matching behavior as minimizing the corresponding Gaussian \(W_2\) term.  My preferred reading is therefore: optimal transport gives the geometric picture of local probability movement, while the DDPM ELBO gives the statistically convenient training objective.

\section{Symmetry, equivariance, and category-like thinking}

If a group $G$ acts on an input space and an output space, a map $F$ is equivariant if
\[
  F(gx)=gF(x).
\]
Convolutional neural networks exploit translation equivariance; graph neural networks exploit permutation equivariance; gauge-equivariant networks try to respect local choices of frame on a bundle.

Here is the matrix calculation behind the slogan.  Let the graph be the cycle \(C_4\), with adjacency matrix
\[
  A=
  \begin{pmatrix}
  0&1&0&1\\
  1&0&1&0\\
  0&1&0&1\\
  1&0&1&0
  \end{pmatrix},
\]
and let \(P\) be the permutation matrix rotating the four vertices by \(90^\circ\).  Since \(P\) is a graph automorphism,
\[
  PA=AP.
\]
The same is true for the usual normalized adjacency \(\widehat A\), because normalization is built from \(A\) and the degree matrix, both invariant under the same rotation:
\[
  P\widehat A=\widehat A P.
\]
A standard graph neural-network layer is
\[
  H^{(l+1)}=\sigma(\widehat A H^{(l)}W),
\]
where \(\sigma\) acts entrywise.  If the input features are rotated, then
\[
\begin{aligned}
  F(PH^{(l)})
  &=\sigma(\widehat A PH^{(l)}W)\\
  &=\sigma(P\widehat A H^{(l)}W)\\
  &=P\sigma(\widehat A H^{(l)}W)\\
  &=PH^{(l+1)}.
\end{aligned}
\]
Equivariance is therefore not only a philosophical commitment to symmetry.  In this example it is the algebraic fact that the permutation matrix commutes with the graph operator.

The original essay also speculates about topos-theoretic or categorical semantics for AI.  A cautious reconstruction is this: category theory is useful when one needs to track structure-preserving maps, functorial changes of representation, and compositional semantics.  It does not by itself explain intelligence.  Its honest role is to provide a disciplined language for invariance, abstraction, and transport of structure.

\section{Further direction}

The note is best read as a map of mathematical languages for data: topology sees holes, geometry sees coordinates and curvature, measure theory sees distributions, dynamics sees evolution, algebra sees symmetry, and category theory sees structure-preserving translation.  The useful question is not which language is fashionable, but which structure is actually stable under the transformations used by the model.

\section{Backpropagation and differential geometry}
Backpropagation is reverse-mode differentiation.  Geometrically, forward computation pushes tangent vectors forward, while backpropagation pulls cotangent vectors backward through the computation graph.
\section{Manifold hypothesis and representation learning}
The manifold hypothesis says high-dimensional data may concentrate near lower-dimensional geometric structures.  Embeddings, charts, curvature, and local neighborhoods then become relevant to learning representations.
\section{Symmetry and equivariance}
If a task has a symmetry group \(G\), an equivariant model satisfies
\[
f(gx)=g f(x).
\]
This reduces sample complexity by building the correct transformation law into the model.
\section{Optimal transport and distributions}
Optimal transport compares probability measures by the cost of moving mass.  It connects geometry, PDE, and machine learning through distances between distributions.
\section{Mathematical structures for learning systems}
Pure mathematics contributes language for structure: cotangent pullback for gradients, manifolds for data geometry, equivariance for symmetry, optimal transport for distributions, and category-like compositionality for architectures.

\section{Rebuilt reading path: what pure mathematics contributes}


There are several layers.

\begin{enumerate}[(i)]
  \item \textbf{Geometry} studies the shape of the representation space: distances, curvature, geodesics, embeddings, and coordinate charts.
  \item \textbf{Topology} studies features stable under continuous deformation: connected components, loops, voids, and persistence across scales.
  \item \textbf{Measure theory} studies data as distributions rather than isolated points.
  \item \textbf{Dynamical systems} study learning or generation as evolution of states or densities.
  \item \textbf{Algebra and category theory} study invariance, symmetry, and composition of structure-preserving maps.
\end{enumerate}

Pure mathematics does not magically explain AI.  It supplies languages in which different kinds of stability can be stated precisely.

\section{Topological data analysis in usable form}

Given a finite metric point cloud $X$, persistent homology constructs a family of complexes $K_\varepsilon$ indexed by a scale parameter.  For the Vietoris--Rips complex,
\[
  \{x_0,\ldots,x_k\}\in VR_\varepsilon(X)
  \quad\Longleftrightarrow\quad
  d(x_i,x_j)\le \varepsilon\text{ for all }i,j.
\]
The homology group is
\[
  H_k(K_\varepsilon)=\ker \partial_k/\operatorname{im}\partial_{k+1}.
\]
This quotient is the mathematical form of the phrase ``holes that are cycles but not boundaries.''

\section{Wasserstein geometry and stochastic dynamics}

For probability measures $\mu,\nu$ on a metric space,
\[
  W_p(\mu,\nu)=\left(\inf_{\pi\in \Pi(\mu,\nu)}\int d(x,y)^p\,d\pi(x,y)\right)^{1/p}.
\]
The coupling $\pi$ is the hidden plan.  In generative modeling, one often tries to construct a transport from a simple distribution to a complicated one.  Normalizing flows do this by invertible maps, diffusion models by stochastic dynamics, and optimal transport by variational minimization.

A diffusion process has the schematic form
\[
  dX_t=b(t,X_t)\,dt+\sigma(t)\,dW_t.
\]
The reverse process depends on the score field $\nabla_x\log p_t(x)$.  Training a score network is therefore an attempt to approximate a time-dependent vector field on probability density.

\section{Category language without exaggeration}

A category consists of objects and morphisms, with associative composition and identity morphisms.  A functor preserves this structure.  Category theory is useful when the same construction can be performed across many spaces in a way compatible with maps.  For example, homology is functorial:
\[
  f:X\to Y \quad\Longrightarrow\quad H_k(f):H_k(X)\to H_k(Y).
\]
The honest categorical contribution is functoriality and compositionality, not a mystical explanation of intelligence.
